\fchapter{Ej 2 Pág 70}

\fsection{\boldit{\textcolor{waveBlue2}{Actividades}}} 
\subsection*{Reflexiona sobre el concepto de pobreza que pueden sufrir algunos colectivos, como las personas mayores, y enumera algunas medidas que pueden adoptar las empresas de distribución de alimentos para contribuir a la consecución del ODS 1.}

\begin{solution}
La pobreza que afecta a colectivos como las personas mayores va más allá de la falta de ingresos. Se trata de una pobreza multidimensional 
que incluye aspectos como la pobreza energética, la dificultad para acceder a una alimentación adecuada, el aislamiento social y las barreras 
físicas o digitales para realizar compras básicas. Este contexto hace que su vulnerabilidad sea mayor y requiera soluciones específicas.

	\begin{enumerate}[label=\alph*)]
		\item\textbf{Descuentos específicos y tarjetas de fidelización solidarias:} 

			Establecer descuentos permanentes para jubilados en días específicos de la semana o crear programas de fidelización donde los puntos acumulados 
			puedan canjearse por productos de primera necesidad o donarse a bancos de alimentos locales.
		

		\item\textbf{Colaboración con servicios sociales y distribución de excedentes:} 

			Crear convenios con ayuntamientos y servicios sociales para identificar a personas mayores en riesgo y facilitarles cupones o tarjetas prepago para la compra. 
			Además, optimizar la logística para donar los excedentes alimentarios diarios (fruta, verdura, panadería) a residencias o centros de día.


		\item\textbf{Servicios de proximidad y accesibilidad:} 

			Ampliar y promocionar servicios como la compra online con entrega a domicilio gratuita o a bajo coste para mayores, 
			y el teléfono de pedido para personas no digitalizadas. También, adaptar las tiendas físicas para mejorar su accesibilidad.
	\end{enumerate}


\end{solution}
